\chapter{Construction of Relative Architecture}\label{chap:1}

\section*{Overview of the Chapter}

In this chapter we construct the objects and maps used to compare the structure of
software. The starting point is to make explicit the choice of what to read as
structure and which operations and laws to preserve.

For example, consider a state consisting of a displayed value and an internal value.
A read returns the displayed value, and an update keeps the internal value and
replaces the displayed value. If a change of the state representation leaves the
displayed value unchanged, reads are preserved. However, applying the change and
then updating need not give the same final state as updating and then applying the
change. To preserve both reads and updates, the two orders must also give the same
result.

Such a comparison requires both the structure of objects and the way operations act
on that structure. In AAT, a typed primitive fact is called an Atom, a family of
Atoms together with relations is called a configuration, and an object that also
keeps state sets, the interpretation of operations, and the like is called an architecture
object. The equations that objects must satisfy are specified as Laws, and we
provide the local structure for reading information part by part and gluing it
together. The choice of vocabulary, extraction rules, object formation, equations,
operations, and locality used in these constructions is called a reading.
Relativizing architecture to a reading means making this choice explicit together
with the objects.

The structure of the chapter and its main results are as follows.

\begin{xltabular}{\linewidth}{@{}LLL@{}}
\toprule
Question & Construction & Outcome \\
\midrule
\endhead
What to take as objects, and which operations to permit & Equip an Atom family with relations, structure, and operations & A core with a family of objects closed under operations (\S\S\ref{sec:1.1}--\ref{sec:1.4}, Theorem~\ref{thm:1.18}) \\
\addlinespace[3pt]
Which information to read locally and glue & Choose contexts, covers, and coefficients & A geometry with a site and a ring for each context (\S\S\ref{sec:1.5}--\ref{sec:1.6}) \\
\addlinespace[3pt]
What to preserve when the reading changes & Define morphisms comparing the data at each stage & The tower of projections connecting extraction, core, and geometry (\S\ref{sec:1.7}, Theorem~\ref{thm:1.31}) \\
\addlinespace[3pt]
How to express existing semantics of CS & Construct objects, operations, and Laws from lenses and protocols & Equivalence of the validity of laws, and recovery of semantics-preserving morphisms (\S\S\ref{sec:1.8}--\ref{sec:1.10}, Propositions~\ref{prop:1.41} and~\ref{prop:1.43}) \\
\bottomrule
\end{xltabular}

\section{Atoms and relative extraction}\label{sec:1.1}

To identify a single fact, one must specify, in addition to its content, what the
fact is about and from which viewpoint it is stated. An Atom carries this
information in five coordinates.

\begin{definition}[Vocabulary of Atoms]\label{def:1.1}

Let $K$, $X$, $S$, $P$, and $T$ be the sets of kinds, axes, subjects, predicates,
and payloads, respectively. The set of Atoms $\mathrm{At}$ is a nonempty set
equipped with maps
\[
\begin{aligned}
\mathrm{kind}&:\mathrm{At}\to K,&
\mathrm{axis}&:\mathrm{At}\to X,&
\mathrm{subject}&:\mathrm{At}\to S,\\
\mathrm{predicate}&:\mathrm{At}\to P,&
\mathrm{payload}&:\mathrm{At}\to T
\end{aligned}
\]
The combined map
$\mathrm{At}\to K\times X\times S\times P\times T$ is required to be injective.
Hence two Atoms are equal if and only if all five coordinates are equal.
The kind expresses the type of the fact, the axis the structure under attention,
and the subject what the fact is about. The predicate and the payload give what is
stated about that subject. For example, the existence of a component, a named
operation, a state quantity of an object, and a relation between two objects can
be expressed by Atoms of different kinds. The payload holds names, types, values,
and the like. Facts treated together are collected into a family of Atoms, and
their relations into a configuration.

\end{definition}

\begin{definition}[Atom families]\label{def:1.2}

An Atom family is a subset $F\subseteq\mathrm{At}$. Its support and its
restriction to an axis $x\in X$ are defined by
\[
\mathrm{supp}(F)=\{\mathrm{subject}(a)\mid a\in F\},
\qquad
F|_x=\{a\in F\mid \mathrm{axis}(a)=x\}
\]
Finiteness of an Atom family means finiteness of $F$.
An extracted family may in general be infinite. The composition reading of
\S\ref{sec:1.2} takes finite families as input.

The data from which extraction is performed are called a source. Even for the same
source, the facts extracted may vary with the choice of vocabulary and resolution.
An extraction doctrine collects this choice together with the conditions for
adopting an Atom.

\end{definition}

\begin{definition}[Extraction doctrines]\label{def:1.3}

An extraction doctrine $D$ consists of a set of sources $\mathrm{Src}_D$,
parameters $v_D,\gamma_D,\rho_D$ for vocabulary, meaning, and resolution, a
normalization map $N_D:\mathrm{Src}_D\to\mathrm{Src}_D$ of sources, and the
following four predicates.

\begin{xltabular}{\linewidth}{@{}LL@{}}
\toprule
Predicate & Meaning \\
\midrule
\endhead
$V_D(v,a)$ & The vocabulary $v$ admits the Atom $a$ \\
\addlinespace[3pt]
$M_D(\gamma,s,a)$ & The meaning reading $\gamma$ admits $a$ on the source $s$ \\
\addlinespace[3pt]
$R_D(\rho,s,a)$ & The resolution $\rho$ admits $a$ on $s$ \\
\addlinespace[3pt]
$E_D(s,a)$ & $a$ holds in the semantics of the source \\
\bottomrule
\end{xltabular}

After normalizing the source, we adopt the Atoms that satisfy all four conditions.
The extraction predicate and the extracted family are defined by
\[
\begin{aligned}
\mathrm{Extracts}_D(s,a)
&\Longleftrightarrow
V_D(v_D,a)\land M_D(\gamma_D,N_D(s),a)\\
&\hspace{2em}\land R_D(\rho_D,N_D(s),a)\land E_D(N_D(s),a),\\
\mathrm{Atomize}_D(s)
&=\{a\in\mathrm{At}\mid \mathrm{Extracts}_D(s,a)\}
\end{aligned}
\]
At this stage we use the structure of $N_D$ as a map; constructions that need
idempotence state it explicitly as an additional condition.

\end{definition}

\begin{proposition}[Existence and uniqueness of the extracted family]\label{prop:1.4}

Fix $D$ and $s\in\mathrm{Src}_D$. Then an Atom family $F$ satisfying
\[
\forall a\in\mathrm{At},\qquad
a\in F\ \Longleftrightarrow\ \mathrm{Extracts}_D(s,a)
\]
exists and is unique.

\end{proposition}

\begin{proof}

The family $\mathrm{Atomize}_D(s)$ of Definition~\ref{def:1.3} satisfies this
condition. If two families $F,G$ satisfy the condition, then
$a\in F\Longleftrightarrow a\in G$ for every Atom, and hence $F=G$ by
extensionality of sets.

\end{proof}

What this uniqueness fixes is the family of facts for a given source and doctrine.
A readout of that family is given by a separate map $o:F\to Y$.
If distinct Atoms satisfy $o(a)=o(b)$, this readout does not distinguish them.
The construction of the extracted family and the amount of information read out of
it can therefore be described separately.

\section{Configurations, objects, and operations}\label{sec:1.2}

An Atom family is the collection of adopted facts. To treat it as a single object,
we add the relations between the facts, together with the state sets and the
interpretation of operations. The relations are kept in a configuration, and the concrete structure
data in an architecture object.

\begin{definition}[Configuration]\label{def:1.5}

A configuration is a triple $C=(F_C,R_C,I_C)$, where $F_C\subseteq\mathrm{At}$
and $R_C,I_C\subseteq\mathrm{At}\times\mathrm{At}$.
Here $R_C$ expresses relations and $I_C$ specified identifications.
When the identifications are used as an equivalence relation, we require its
reflexivity, symmetry, and transitivity.
When $R_C,I_C\subseteq F_C\times F_C$, we say that the relations and
identifications are supported by the family.

A composition reading is a rule that assigns to each finite family $F$ a
configuration
\[
\mathrm{Comp}(F)=(F,R_F,I_F),
\qquad R_F,I_F\subseteq F\times F
\]
A configuration obtained by restricting the relations and identifications to a
finite subfamily is called a molecule. It is the finite unit for handling several
Atoms together with their relations when comparing local structure.

\end{definition}

\begin{definition}[Architecture object]\label{def:1.6}

An architecture object is a triple $A=(C_A,S_A,Q_A)$, where $C_A$ is a
configuration, $S_A$ is structure data of a chosen type, and $Q_A$ is a quantity
of a chosen type; the types themselves are included in the data of the object.
The structure data may hold state sets and transition maps, graphs, diagrams,
algebras, and the like. The quantity may hold specified values or results of
evaluation. The totality of such objects, in a universe of sufficient size, is
written $\mathrm{ArchObj}(\mathrm{At})$.

An object reading is a rule that sends a configuration $C$ to an object
$\mathrm{Form}(C)$ with $C_{\mathrm{Form}(C)}=C$.
Fixing this rule, from a finite family we obtain the construction
\[
F\longmapsto \mathrm{Comp}(F)\longmapsto
\mathrm{Form}(\mathrm{Comp}(F))
\]
One may also change the state sets or the interpretation of operations while
keeping the configuration the same, in which case the result is a different
architecture object.

\end{definition}

\begin{definition}[Morphisms of configurations and operations]\label{def:1.7}

An operation between objects is accompanied by an action on Atoms, relations, and
identifications. A morphism of configurations expresses this action.

A morphism of configurations $h:C\to D$ is a map
$h_{\mathrm{At}}:\mathrm{At}\to\mathrm{At}$ such that
\[
\begin{aligned}
a\in F_C&\Longrightarrow h_{\mathrm{At}}(a)\in F_D,\\
(a,b)\in R_C&\Longrightarrow
(h_{\mathrm{At}}(a),h_{\mathrm{At}}(b))\in R_D,\\
(a,b)\in I_C&\Longrightarrow
(h_{\mathrm{At}}(a),h_{\mathrm{At}}(b))\in I_D
\end{aligned}
\]
With identity maps and composition of maps, configurations form a category.
Indeed, the three implications all hold for the identity map and are closed under
composition.

An operation reading consists of a set $\mathrm{Op}(A,B)$ for each pair of
objects and maps
\[
d_{A,B}:\mathrm{Op}(A,B)\longrightarrow
\mathrm{Hom}(C_A,C_B)
\]
An element $o\in\mathrm{Op}(A,B)$ is called an operation from $A$ to $B$.
An operation keeps its own name and additional data, and $d_{A,B}$ gives its
action on configurations. Even when two operations have the same action, the
operations themselves are kept distinct.
If identity operations and composition are specified so that the unit and
associativity laws hold and $d$ preserves them, the objects and operations form a
category.

\end{definition}

\begin{example}[Named operations with the same action]\label{ex:1.8}

Let $e_1,e_2$ be two operations on the state set $\{0,1\}$, both acting as the
identity map. Taking the operation set to be $\{e_1,e_2\}$, one can distinguish a
condition specified for $e_1$ from a condition specified for $e_2$ even though
their execution results are equal.
In \S\ref{sec:1.10} we actually construct morphisms of configurations whose Atoms
include the operation names.

\end{example}

\begin{definition}[Invariants and signatures]\label{def:1.9}

Values and properties to be preserved by operations are specified as invariants,
and the tuple of values used for comparison is collected into a signature.

A functional invariant is a map $I:\mathrm{ArchObj}(\mathrm{At})\to V_I$, and a
predicate invariant is a predicate $P(A)$ on objects.
Preservation under an operation $o:A\to B$ means
$I(A)=I(B)$ and $P(A)\Rightarrow P(B)$, respectively.
For ordered quantities, the condition $I(B)\leq I(A)$ can also be specified
separately. An invariant reading chooses an index set and assigns to each index a
functional or predicate invariant.

A signature reading consists of a set of axes $\Lambda$, a value range
$V_\lambda$ for each axis, a subset $\Lambda_{\mathrm{sel}}$ of selected axes,
and values $q_\lambda(A)\in V_\lambda$ for each object.
The signature of an object is
$(q_\lambda(A))_{\lambda\in\Lambda_{\mathrm{sel}}}$.
Selecting finitely many axes yields a finite multi-axis representation.

\end{definition}

\section{Local contexts and Laws}\label{sec:1.3}

Instead of handling all the information of an object at once, we consider
extracting part of its structure and observed values. A context specifies at
which positions, along which axes, and what can be read. Expressing an equation
on a context as a residual, so that the equation holds when the value is zero, allows Laws to be
evaluated locally.

\begin{definition}[Contexts and restrictions]\label{def:1.10}

A context $W$ of an object $A$ carries three sets
$\mathrm{Supp}(W)$, $\mathrm{Ax}(W)$, and $\mathrm{Obs}(W)$, together with
\[
\mathrm{reads}_W\subseteq\mathrm{Supp}(W)\times\mathrm{At},
\quad
\mathrm{readAx}_W\subseteq\mathrm{Ax}(W),
\quad
\mathrm{readObs}_W\subseteq\mathrm{Obs}(W)
\]
An element of $\mathrm{Obs}(W)$ is called an observable.
If $\mathrm{reads}_W(s,a)$ holds, we require $a\in F_{C_A}$.
An element of the support is a position at which Atoms are read; its role differs
from that of the set of subjects in Definition~\ref{def:1.2}.
Any additional data needed are also included in the context.

A context reading chooses a small preorder of contexts.
The relation $W'\leq W$ expresses the direction in which data obtained on $W$ can
be restricted to $W'$.
For each $W'\leq W$ we specify maps
\[
\mathrm{Supp}(W')\longrightarrow\mathrm{Supp}(W),\qquad
\mathrm{Ax}(W')\longrightarrow\mathrm{Ax}(W),\qquad
\mathrm{Obs}(W)\longrightarrow\mathrm{Obs}(W')
\]
The support map preserves the reading of the same Atom, the axis map preserves
readable axes, and the restriction of observables sends readable observables to
readable ones.
The three maps are compatible with identities and composition in their respective directions.
We regard this preorder as a category and write $W'\to W$ for the morphism
corresponding to $W'\leq W$.
The context category is denoted $\mathcal{C}_A$.

\end{definition}

\begin{example}[Contexts given by subsets]\label{ex:1.11}

For a finite family $F$, take the subsets $U\subseteq F$ as contexts and the
inclusions as morphisms. Setting $\mathrm{Supp}(U)=U$, with each element read as
its own Atom, taking the axes to be a one-element set, and taking
$\mathrm{Obs}(U)$ to be the set $B^U$ of functions into a set $B$, the
restriction of observables is the restriction of functions.
The intersection of two contexts is given by $U\cap V$.

To handle equations, we choose a ring in which local values can be added and
subtracted. Taking the difference of the two sides of an equation expresses its
validity as the vanishing of a residual. Below we keep both the symbolic
coordinates that present the equations and the residuals evaluated on objects.

\end{example}

\begin{definition}[Architectural equation system]\label{def:1.12}

Fix a context category $\mathcal{C}$ and a set of architecture objects.
An equation system $E$ consists of the following data.

\begin{itemize}
\item A set $K_E$ of equation indices and a role for each index,
\[
\mathrm{role}_E:K_E\to\{\mathrm{required},\mathrm{optional},\mathrm{derived}\}.
\]
\item A presheaf of commutative rings
$O_E:\mathcal{C}^{\mathrm{op}}\to{\mathbf{CommRing}}$.
The restriction along a morphism $j:W'\to W$ is written
$\mathrm{res}_j:O_E(W)\to O_E(W')$.
\item A symbolic coordinate $\nu_{W,i,a}\in O_E(W)$ for each context, index, and
Atom.
\item A residual $\varepsilon_{W,A,i,a}\in O_E(W)$ depending in addition on an
object $A$.
\end{itemize}

By a presheaf we mean here that the rings and restriction homomorphisms satisfy
$\mathrm{res}_{\mathrm{id}}=\mathrm{id}$ and
$\mathrm{res}_{jk}=\mathrm{res}_k\mathrm{res}_j$.
On the two families of coordinates we impose
\begin{equation*}
\mathrm{res}_j(\nu_{W,i,a})=\nu_{W',i,a},
\qquad
\mathrm{res}_j(\varepsilon_{W,A,i,a})=\varepsilon_{W',A,i,a}
\tag{1.1}\label{eq:1.1}
\end{equation*}
The $\nu$ are coordinates for presenting the equations algebraically, and the
$\varepsilon$ are residuals evaluating whether the equations hold on an object.
The Law $L_i$ is the name for the equation with index $i$ together with its
meaning.

\end{definition}

\begin{definition}[Lawfulness]\label{def:1.13}

The validity of an equation and the condition of satisfying all required
equations are defined by
\[
\begin{aligned}
E_i(A)
&\Longleftrightarrow
\forall W\in\mathrm{Ob}(\mathcal{C})\ \forall a\in\mathrm{At},\
\varepsilon_{W,A,i,a}=0,\\
\mathrm{Lawful}_E(A)
&\Longleftrightarrow
\forall i\in K_E^{\mathrm{req}},\ E_i(A),\\
K_E^{\mathrm{req}}
&=\{i\in K_E\mid\mathrm{role}_E(i)=\mathrm{required}\}
\end{aligned}
\]
That is, lawfulness means that for each required equation the residual vanishes
on every context and Atom. Validity at all indices is written
$\mathrm{FullyLawful}_E(A)$.
The roles required, optional, and derived designate an equation as mandatory, as
optional, and as one to be derived, respectively.
To conclude the validity of a derived equation from the other equations, one
proves the implication.

\end{definition}

\section{Finite detection and generation of the core}\label{sec:1.4}

The failure of an equation can sometimes be detected by inspecting a small number
of Atoms and relations. To describe this finite detection, we define queries
about structure and their combinations. We then collect objects, operations,
equations, and detection methods, and construct as the core the totality of
objects reachable by operations.

\begin{definition}[Signed finite queries]\label{def:1.14}

A query on a configuration is one of the following: membership of an Atom $a$,
validity of a relation $R(a,b)$, or validity of an identification $I(a,b)$.
Queries for relations and identifications also require that the Atoms at both
ends belong to the family.
We write $\mathrm{Holds}_A(q)$ when a query $q$ holds on an object $A$.

A signed finite query sequence is $Q=((q_1,b_1),\ldots,(q_n,b_n))$ with
$b_j\in\{0,1\}$, and we define
\[
\mathrm{Matches}(Q,A)
\Longleftrightarrow
\forall j,\quad \mathrm{Holds}_A(q_j)\Longleftrightarrow b_j=1
\]
The sign $0$ allows the absence of an Atom or a relation to be expressed as a
finite condition as well.
For example, attaching the signs $1,1,0$ to the three queries ``$a$ belongs'',
``$b$ belongs'', and ``$R(a,b)$ holds'' expresses the pattern in which the two
Atoms exist but the specified relation between them does not.

\end{definition}

\begin{definition}[Finite detectors and circuits]\label{def:1.15}

A detector code is obtained by finitely many applications of the constructors
$\mathrm{reject}$, which expresses the empty selection, $\mathrm{exact}(Q)$,
which specifies a single sequence, and $\mathrm{any}(c,d)$, the disjunction of
two codes.
The sequences it accepts are, respectively, the empty set, $\{Q\}$, and the union
of the accepted sets.
To each $i\in K_E$ we assign a code $c_i$ and write $T_i$ for its finite
accepted set.

The set of circuits on an object $A$ is
\[
\mathrm{Circ}_E(A,i)=\{Q\in T_i\mid\mathrm{Matches}(Q,A)\}
\]
That is, a circuit is a finite pattern, among those specified by the code, that actually holds
on the object.
To use circuits as evidence for the failure of an equation, we require the
soundness below.
The condition that the failure of every required equation can be detected is
called completeness.
Each condition is required for all objects $A$, indices $i$, and query
sequences $Q$.
\[
\begin{aligned}
Q\in\mathrm{Circ}_E(A,i)&\Longrightarrow \neg E_i(A),\\
i\in K_E^{\mathrm{req}}\ \land\ \neg E_i(A)
&\Longrightarrow\mathrm{Circ}_E(A,i)\ne\varnothing.
\end{aligned}
\]
Soundness states that the failure of the equation follows from an accepted
structural pattern.
Completeness states that a failure of a required equation always has such a
finite pattern.
Both are conditions to be proved for the specified equation system and detector.

\end{definition}

\begin{proposition}[Deciding lawfulness by circuits]\label{prop:1.16}

If the detector is sound, all required circuits of a lawful object are empty.
If it is moreover complete for the required equations, then
\[
\mathrm{Lawful}_E(A)
\Longleftrightarrow
\forall i\in K_E^{\mathrm{req}},\
\mathrm{Circ}_E(A,i)=\varnothing.
\]

\end{proposition}

\begin{proof}

If a lawful object had a required circuit, soundness would make the
corresponding required equation fail, a contradiction.
Conversely, if some required equation fails, completeness provides a
corresponding circuit.

\end{proof}

The same applies when the presence or absence of failures is expressed by quantities
$\omega_i(A)$. Given $E_i(A)\Longleftrightarrow\omega_i(A)=0$ and an aggregation
with $\mathrm{Agg}((v_i)_i)=0\Longleftrightarrow\forall i,\ v_i=0$, the vanishing
of the aggregation over the required indices is equivalent to lawfulness.
This equivalence follows by composing the two equivalences.

\begin{definition}[Core reading]\label{def:1.17}

A core reading $r$ includes the following choices.

\begin{itemize}
\item \textbf{Extraction}: choose a doctrine and a source $(D,s)$, and require
the extracted family $F_r=\mathrm{Atomize}_D(s)$ to be finite.
\item \textbf{Object formation}: choose a composition reading and an object
reading.
\item \textbf{Operations and values for comparison}: choose an operation
reading, an invariant reading, and a signature reading.
\end{itemize}

From these we define the initial configuration and object
\[
C_r=\mathrm{Comp}_r(F_r),\qquad A_r=\mathrm{Form}_r(C_r)
\]
We further choose the context category $\mathcal{C}_r$ of the base object $A_r$,
an equation system $E_r$ on it, and a sound detector.
The residuals are required to be evaluable on every object of
$\mathrm{ArchObj}(\mathrm{At})$, and the context category chosen from the base
object is used in common for evaluating the generated family of objects.

A subset $Z\subseteq\mathrm{ArchObj}(\mathrm{At})$ of objects is closed under
operations if $A\in Z$ and $o\in\mathrm{Op}_r(A,B)$ imply $B\in Z$.
Let $\mathrm{Obj}_r$ be the intersection of all such subsets containing the base
object.
The core $\mathrm{Core}(r)$ is the combined data of this family of objects, the
base object, the context category, the equation system, the family of circuits,
the family of operations, and the invariants and signature.

\end{definition}

\begin{theorem}[Generation of the relative core]\label{thm:1.18}

From a core reading $r$ the core $\mathrm{Core}(r)$ is constructed.
The family of its initial configuration equals $F_r$, and the configuration of
the base object equals $C_r$.
Moreover, $\mathrm{Obj}_r$ is the smallest family of objects containing the base
object and closed under operations, and it is also characterized by the
following finite reachability.
\begin{equation*}
B\in\mathrm{Obj}_r
\Longleftrightarrow
\exists n\in\mathbb{N},\
A_r=A_0\xrightarrow{o_1}A_1\xrightarrow{o_2}
\cdots\xrightarrow{o_n}A_n=B.
\tag{1.2}\label{eq:1.2}
\end{equation*}
Sequences of length $0$ are allowed.

\end{theorem}

\begin{proof}

The extracted family is obtained by Proposition~\ref{prop:1.4}, and the
finiteness condition of the reading allows it to be fed into the composition.
The conditions on the composition and object readings yield the two equalities
of configuration components.
Since the set of all objects contains the base object and is closed under
operations, the family over which the intersection of
Definition~\ref{def:1.17} is taken is nonempty.

The intersection contains the base object. If $A$ belongs to the intersection
and there is an operation $o:A\to B$, then $B$ belongs to every subset defining
the intersection, so the intersection is also closed under operations.
Minimality follows from the definition of the intersection.

Let $Z_{\mathrm{fin}}$ be the set of objects reachable from the base object by
finite sequences.
By the sequences of length $0$ and the appending of an operation to the end of a
sequence, $Z_{\mathrm{fin}}$ contains the base object and is closed under
operations.
Hence $\mathrm{Obj}_r\subseteq Z_{\mathrm{fin}}$.
The reverse inclusion is obtained by induction on the length of the sequence.
Finally, restricting the specified equations, detector, invariants, and
signature to this family of objects and taking the operations for each pair of
objects gives all the components of the core.

\end{proof}

The finiteness in \eqref{eq:1.2} is a condition on the length of each witness of
reachability.
The size of the whole family of objects closed under operations is determined by
the specified operation reading.

\section{From covers to sites}\label{sec:1.5}

To read structure part by part, one needs a family of contexts covering the
whole and a place where the information read on two parts can be compared.
For the subsets of Example~\ref{ex:1.11}, the latter is the intersection.
For general contexts we choose this overlap as a pullback, and take as the
starting point of covers the families that can jointly read the required
information among Atoms, equations, and signatures.

\begin{definition}[Overlaps and coverage requirements]\label{def:1.19}

For each diagram $U\to W\leftarrow V$ of two morphisms in the context category
$\mathcal{C}$, we choose an object $U\times_W V$ and projections $\pi_U,\pi_V$.
The composites of the two projections with $U\to W$ and $V\to W$, respectively,
are required to be equal.
Moreover, any $T\to U$ and $T\to V$ satisfying this commutativity are obtained
by composing a unique $T\to U\times_W V$ with the projections.
This pullback is called the overlap, and its choice is written $\mathrm{Ov}$.
When the preorder of contexts is a partial order, the overlap is the greatest
lower bound of the two contexts below $W$.

Coverage requirements $\mathcal{R}$, with an equation system $E$ and a signature
fixed, choose the following.

\begin{itemize}
\item A subset $A_{\mathrm{req}}\subseteq\mathrm{At}$ of required Atoms.
\item Required equation coordinates
$C_{\mathrm{req}}\subseteq K_E^{\mathrm{req}}\times\mathrm{At}$.
\item A subset $C_{\mathrm{wit}}\subseteq K_E\times\mathrm{At}$ of indices of
the symbolic coordinates $\nu_{W,i,a}$ to be read as witnesses (evidence of
violation).
\item Required signature axes $\Lambda_{\mathrm{req}}\subseteq\Lambda$.
\item Four visibility predicates expressing that these can be read on each
context, and a predicate $B_{\mathcal{R}}(U,W)$ expressing that the required
interaction can be read on the overlap.
\end{itemize}

The selected coordinates and witness indices are kept along the specified
restrictions.
For example, the restriction \eqref{eq:1.1} restricts the value while keeping
the index $(i,a)$.
Visibility is a separate condition for a context to actually read the data of
that index.

A family of morphisms $\mathcal{U}=(u_\alpha:W_\alpha\to W)_{\alpha\in I}$ is
$(E,\mathcal{R},\mathrm{Ov})$-admissible if it satisfies the following
conditions.

\begin{itemize}
\item \textbf{Information read on the parts}: each required Atom and signature
axis is readable on at least one $W_\alpha$.
\item \textbf{Information read on parts or overlaps}: each required equation
coordinate and selected witness is readable on at least one $W_\alpha$ or
$W_{\alpha\beta}=W_\alpha\times_W W_\beta$.
\item \textbf{Interaction on overlaps}: $B_{\mathcal{R}}(W_{\alpha\beta},W)$
holds for all $\alpha,\beta$.
\item \textbf{Correspondence with the object}: the Atoms read by the support map
of each leg belong to the family of the underlying object.
\end{itemize}

A system of covers is required to be closed under restriction to another context
and under further covering of each part.
We define this system using the following sieves.

\end{definition}

\begin{definition}[Sieves and Grothendieck topologies]\label{def:1.20}

A sieve $S$ on $W$ is a collection of morphisms with codomain $W$ such that if
$f:V\to W$ belongs to it, then so does $fg$ for every $g:V'\to V$.
The pullback along a morphism $u:W'\to W$ is
$u^*S=\{g:V\to W'\mid ug\in S\}$.
The sieve $\langle \mathcal{U}\rangle$ generated by a family $\mathcal{U}$
consists of all morphisms that factor through at least one $u_\alpha$.

A Grothendieck topology $J$ assigns to each $W$ a set $J(W)$ of covering sieves
such that the following hold.

\begin{enumerate}
\item The maximal sieve consisting of all morphisms with codomain $W$ belongs
to $J(W)$.
\item If $S\in J(W)$, then $u^*S\in J(W')$.
\item If $S\in J(W)$ and another sieve $T$ satisfies $u^*T\in J(W')$ for every
$u:W'\to W$ with $u\in S$, then $T\in J(W)$.
\end{enumerate}

We use the definition of
\cite[\href{https://stacks.math.columbia.edu/tag/00YW}{\S7.47, Tag 00YW}]{Stacks}.

\end{definition}

\begin{proposition}[Generation of the topology from coverage requirements]\label{prop:1.21}

There exists a smallest Grothendieck topology containing the sieves generated by
all admissible families.
We write it $J_{E,\mathcal{R},\mathrm{Ov}}$.

\end{proposition}

\begin{proof}

The topology in which every sieve on every object is a cover contains all the
specified sieves.
Take, objectwise, the intersection of all topologies containing them.
The maximal-sieve condition and stability under pullback hold in each topology
and hence in the intersection.
As for transitivity, if its hypothesis holds in the intersection, it holds in
each topology, and the concluding sieve belongs to all of them.
The intersection is therefore a topology, and it is smallest by definition.

\end{proof}

A site is a pair of a small category and a Grothendieck topology.
The pair $(\mathcal{C}_r,J_{E_r,\mathcal{R},\mathrm{Ov}})$ obtained from the
base object $A_r$ is called the AAT site.
The choices of the equation system, the signature, the coverage requirements,
and the overlaps are also kept as its construction data.
A covering family of the generated topology that also satisfies the specified
visibility and interaction requirements is called
$(E,\mathcal{R},\mathrm{Ov})$-adequate.
The topology gives the rules for local gluing, and adequacy expresses that the
information used in the theorems is available on the cover.
It is used in Chapter~\ref{chap:2} to connect the local validity of equations
with the consistency of gluing.

\section{Presheaves, sheaves, and geometries with coefficients}\label{sec:1.6}

A presheaf assigns data to each context and has rules for restricting them to
smaller contexts.
A sheaf satisfies in addition the condition that from local data agreeing on the
overlaps, the data on the whole are obtained uniquely.

\begin{definition}[Presheaves and sheaves]\label{def:1.22}

A set-valued presheaf is a contravariant functor
$F:\mathcal{C}^{\mathrm{op}}\to{\mathbf{Set}}$.
An element of $F(W)$ is called a section over $W$.
A matching family on a family $\mathcal{U}=(u_\alpha:W_\alpha\to W)_\alpha$ is a
family of elements $s_\alpha\in F(W_\alpha)$ satisfying
\begin{equation*}
F(\pi_\alpha)(s_\alpha)=F(\pi_\beta)(s_\beta)
\quad\text{in }F(W_{\alpha\beta})
\tag{1.3}\label{eq:1.3}
\end{equation*}
for all $\alpha,\beta$.
The set of matching families is written $\mathrm{Match}(\mathcal{U},F)$.
A presheaf $F$ is a sheaf for $J$ if, for every family with
$\langle \mathcal{U}\rangle\in J(W)$, the map
\begin{equation*}
F(W)\longrightarrow\mathrm{Match}(\mathcal{U},F),
\qquad s\longmapsto(F(u_\alpha)(s))_\alpha
\tag{1.4}\label{eq:1.4}
\end{equation*}
is a bijection. Injectivity is called separatedness, and surjectivity the
existence of gluings.
A morphism of sheaves is a natural transformation of presheaves, and the
category of sheaves is written $\mathrm{Sh}(\mathcal{C},J)$.

A matching family on a sieve $S$ assigns to each $u:V\to W$ with $u\in S$ an
element $s_u\in F(V)$ satisfying $s_{uv}=F(v)(s_u)$.
A matching family on a covering family becomes, by restriction to the morphisms
factoring through the $u_\alpha$, a matching family on the generated sieve.
The values obtained from two factorizations are equal by the universal property
of the pullback and \eqref{eq:1.3}.
Conversely, reading a family on the sieve at each $u_\alpha$ yields
\eqref{eq:1.3}.
The definition therefore also agrees with the definition of sheaves by covering
sieves.

As an example in which the sheaf condition is needed, take the subsets of the
two-point set $\{p,q\}$ as contexts and the families covering by unions as
covers. Consider the presheaf $F$ that takes a one-element set on the empty
context and $\{0,1\}$ on the nonempty contexts, with the identity map as the
restriction between nonempty contexts.
The values $s_p=0$ and $s_q=1$ on the one-point contexts agree on the overlap,
but they glue to neither element over the two points.
Sheafification is the operation that fills such gaps by adding compatible local
data as new sections.

\end{definition}

\begin{proposition}[Sheafification]\label{prop:1.23}

For every presheaf $F$ there exist a sheaf $a_JF$ and a natural transformation
$\eta_F:F\to a_JF$. For every sheaf $G$ and natural transformation
$t:F\to G$ there is a unique $\bar t:a_JF\to G$ with $t=\bar t\eta_F$.

\end{proposition}

\begin{proof}[Construction and proof]

We treat a matching family on a cover as a single section and identify
presentations that agree locally.
Performing this operation twice yields a sheaf equipped with all the required
gluings.

An element of $F^+(W)$ is presented by a matching family on a covering sieve
$S\in J(W)$.
Two presentations are equivalent when they agree on a covering sieve contained
in both.
A finite intersection of covering sieves is again a cover; this follows by
applying the transitivity of Definition~\ref{def:1.20} to the pullbacks along
the morphisms of one of the sieves.
The equivalence relation is therefore transitive.
Pullback defines the restrictions of $F^+$, and the families of restrictions of
global sections give $F\to F^+$.

If two elements of $F^+$ are equal on a cover, we can collect the local covering
sieves expressing this equality. On the cover obtained by the transitivity of
the topology the presentations agree, so the two elements are equal.
Hence $F^+$ is separated.
When $F$ itself is separated, each element of a matching family of $F^+$ is
presented by local sections of $F$. Equality on the overlaps holds after further
localization, and by the separatedness of $F$ it holds on the overlaps themselves.
The family on a covering sieve collecting these presentations gives the gluing.
Thus $F^+$ is a sheaf whenever $F$ is separated.

We therefore set $a_JF=(F^+)^+$.
A transformation $t:F\to G$ sends each presentation to a matching family of $G$
and, by the unique gluing in $G$, determines $F^+\to G$.
Performing the same operation twice gives $\bar t$.
Each element is determined by a local presentation, and since $G$ is separated,
the extension is unique.

\end{proof}

This is the standard sheafification
(\cite[\href{https://stacks.math.columbia.edu/tag/00ZG}{\S7.49, Tag 00ZG}]{Stacks}).
In particular, if the original $F$ is a sheaf, then $\eta_F$ is an isomorphism
by the universal property.

\begin{example}[Sheafification adding local sections]\label{ex:1.24}

The sheafification of the presheaf $F$ on the two-point set above is
$U\mapsto\{0,1\}^U$.
A family taking a different value at each point is also retained, as a single
function.
Indeed, a function is determined uniquely by its restrictions to the single
points, and arbitrary values at the single points can be glued.
From the original presheaf, each value is sent to the constant function.
A natural transformation to any sheaf is determined by its values on the
one-point contexts, and gluing these extends it uniquely from
$U\mapsto\{0,1\}^U$.
The universal property of Proposition~\ref{prop:1.23} is therefore satisfied as
well.

To treat equations algebraically, we give the local data the structure of a
ring.
Taking the coordinates of each context as variables and imposing the structural
relations as polynomial equations produces the ring used on that context.
Specifying restrictions between contexts then yields a ring-valued presheaf.

\end{example}

\begin{definition}[Rings for each context and restrictions]\label{def:1.25}

Choose a coefficient ring $k$. To each context $W$ we assign a set of
coordinates $Z_W$, the kind of each coordinate and the type of its local data,
an index set $H_W$ of structural relations, and polynomials
$r_{W,h}\in k[x_z\mid z\in Z_W]$.
Each element of the polynomial ring is a finite sum in finitely many variables.
The ideal generated by the structural relations and the quotient are
\[
I_W^{\mathrm{str}}=(r_{W,h}\mid h\in H_W),
\qquad
B_{\mathrm{raw}}(W)=k[x_z\mid z\in Z_W]/I_W^{\mathrm{str}}
\]
To each morphism $j:W'\to W$ we assign a $k$-algebra homomorphism
$\rho_j:k[x_z\mid z\in Z_W]\to k[x_z\mid z\in Z_{W'}]$ sending variables to
polynomials, and require
\begin{equation*}
\rho_j(I_W^{\mathrm{str}})\subseteq I_{W'}^{\mathrm{str}},
\qquad
\rho_{\mathrm{id}}=\mathrm{id},
\qquad
\rho_{jk}=\rho_k\rho_j
\tag{1.5}\label{eq:1.5}
\end{equation*}
The coordinates, the relations, and these restrictions of polynomials together
are called a raw restriction system $\mathcal{B}$.
By \eqref{eq:1.5} the restrictions descend to the quotients, and
$B_{\mathrm{raw}}$ becomes a presheaf of commutative $k$-algebras.
This claim follows because the difference of two representatives of the same
residue class stays in the structural ideal after restriction, and because
identities and composition hold at the level of polynomials.

An architecture geometry with coefficients is
\[
G=(r,\mathcal{R},\mathrm{Ov},k,\mathcal{B})
\]
The base object, the contexts, the equations, and the signature come from $r$,
and the topology from Proposition~\ref{prop:1.21}.
The structural relations determine, on each context, the ring representing the
local data.
Chapter~\ref{chap:2} gives a realization connecting this presheaf of rings with
the equation system, and makes explicit the condition under which the
evaluation of the symbolic coordinates corresponds to the residuals of each
object.

\end{definition}

\begin{lemma}[Change of coefficients for raw systems]\label{lem:1.26}

A ring homomorphism $\varphi:k\to k'$ yields a raw system
$\varphi_!\mathcal{B}$ by applying $\varphi$ to the coefficients of all
polynomials while keeping the indices of coordinates and relations.
This construction respects identity homomorphisms and composition of
homomorphisms.

\end{lemma}

\begin{proof}

Send each variable $x_z$ to the same variable and each coefficient $c$ to
$\varphi(c)$.
The new structural ideal is the ideal generated by the images of the structural
relations.
That the original restrictions preserve the structural ideal is expressed by
equations presenting the image of each generating relation as a finite ideal
combination, and these equations remain valid after applying the coefficient
map.
The same applies to the identity and composition equations of \eqref{eq:1.5}.
Successive application of coefficient maps equals application of their
composite, which gives the last claim.

\end{proof}

\section{Morphisms of readings and the three-level projections}\label{sec:1.7}

Changing the reading changes the extracted Atoms, the way objects are formed,
and the choice of equations and locality.
To compare the situations before and after a change, one needs maps connecting
the respective data and the conditions they satisfy.
We define three categories according to the range of data compared.

\begin{xltabular}{\linewidth}{@{}LLL@{}}
\toprule
Category & Data carried by an object & What a morphism compares \\
\midrule
\endhead
$B$ & Extraction doctrine and source & Sources and Atoms \\
\addlinespace[3pt]
$E_{\mathrm{core}}$ & Core reading & In addition to extraction: object formation, operations, equations, detection, invariants, and signatures \\
\addlinespace[3pt]
$E_{\mathrm{geom}}$ & Geometry with coefficients & In addition to the core: covers, overlaps, coefficients, presentations of rings, and the data of contexts \\
\bottomrule
\end{xltabular}

We require the validity of extraction and the validity of equations to be
equivalent before and after the change.
A change equipped with this preservation and reflection is called an exact
change.

\begin{definition}[The category of extraction bases]\label{def:1.27}

Fix the vocabulary of Atoms. An object of the category $B$ is a pair $(D,s)$ of
an extraction doctrine and a selected source. A morphism $(D,s)\to(D',s')$ is a
map $f:\mathrm{Src}_D\to\mathrm{Src}_{D'}$ together with a bijection
$e:\mathrm{At}\xrightarrow{\sim}\mathrm{At}$ satisfying
\begin{equation*}
\begin{aligned}
f(s)&=s',&
N_{D'}f&=fN_D,\\
\mathrm{Extracts}_D(t,a)
&\Longleftrightarrow
\mathrm{Extracts}_{D'}(f(t),e(a))
&&\text{for all }t,a
\end{aligned}
\tag{1.6}\label{eq:1.6}
\end{equation*}
The identity morphism consists of the two identity maps, and composition is
$(f',e')(f,e)=(f'f,e'e)$.
Composing the equalities and equivalences of \eqref{eq:1.6} shows that the same
conditions hold after composition.

The direct image of a family is written $e_*F=\{e(a)\mid a\in F\}$.
From \eqref{eq:1.6} and the bijectivity of $e$ we obtain
\begin{equation*}
\mathrm{Atomize}_{D'}(f(t))=e_*\mathrm{Atomize}_D(t)
\tag{1.7}\label{eq:1.7}
\end{equation*}
Indeed, an element of the right-hand side belongs to the left-hand side by
preservation of extraction.
For an element $b$ of the left-hand side, take $a=e^{-1}(b)$ and use reflection
of extraction.
The source map $f$ is a general map; for example, it may send distinct states
to a single one.

The direct image of a configuration is defined by letting $e$ act on the family
and on both relations.
On finite query sequences and detector codes, $e$ acts on all the Atoms that
occur. These are also written $e_*$.

In the following conditions, the maps are chosen so that, even when the
presentation of objects changes, the validity or failure of the Laws and the
finite patterns witnessing failures correspond.

\end{definition}

\begin{definition}[Morphisms of cores]\label{def:1.28}

The objects of $E_{\mathrm{core}}$ are the core readings with a fixed vocabulary
of Atoms.
A morphism from $r$ to $r'$ consists of a base morphism $(f,e)$ and the
following data and conditions.

\textbf{Object formation and operations.}
We require that the route mapping the Atoms first and then forming the object
agree with the route forming the object first and then mapping it.
We give a map $H:\mathrm{ArchObj}(\mathrm{At})\to\mathrm{ArchObj}(\mathrm{At})$
and require, for every finite family $F$, configuration $C$, and object $A$,
\begin{equation*}
\begin{aligned}
\mathrm{Comp}_{r'}(e_*F)&=e_*\mathrm{Comp}_r(F),\\
H(\mathrm{Form}_r(C))&=\mathrm{Form}_{r'}(e_*C),\\
C_{H(A)}&=e_*C_A
\end{aligned}
\tag{1.8}\label{eq:1.8}
\end{equation*}
On the operation maps
$\Phi_{A,B}:\mathrm{Op}_r(A,B)\to\mathrm{Op}_{r'}(H(A),H(B))$ we impose, for
the action $d(o)$ on configurations,
\begin{equation*}
d(\Phi(o))_{\mathrm{At}}\,e=e\,d(o)_{\mathrm{At}}
\tag{1.9}\label{eq:1.9}
\end{equation*}
Since $e$ itself gives a morphism of configurations $C_A\to C_{H(A)}$,
\eqref{eq:1.9} is also a commutative square of configurations.

\textbf{Equations and detection.}
We match up the contexts, the equation indices, and the rings in which values
are taken, and preserve the coordinates, residuals, and detector codes under
this correspondence.
We give an equivalence of categories
$\theta:\mathcal{C}_r\simeq\mathcal{C}_{r'}$ with a specified quasi-inverse, a
bijection of indices $\lambda:K_{E_r}\xrightarrow{\sim}K_{E_{r'}}$, and ring
isomorphisms $\alpha_W:O_{E_r}(W)\xrightarrow{\sim}O_{E_{r'}}(\theta W)$ for
each context. These are required to satisfy
\begin{equation*}
\begin{aligned}
\mathrm{role}_{E_{r'}}(\lambda i)&=\mathrm{role}_{E_r}(i),\\
\alpha_{W'}\mathrm{res}_j&=\mathrm{res}'_{\theta j}\alpha_W,\\
\alpha_W(\nu_{W,i,a})&=\nu'_{\theta W,\lambda i,e(a)},\\
\alpha_W(\varepsilon_{W,A,i,a})
&=\varepsilon'_{\theta W,H(A),\lambda i,e(a)},\\
c'_{\lambda i}&=e_*c_i
\end{aligned}
\tag{1.10}\label{eq:1.10}
\end{equation*}
The last equality is an equality of detector codes as syntax.

\textbf{Invariants and signatures.}
We give an index map $\mu$ preserving the kind of invariant.
For corresponding functional invariants, we require a bijection $t_i$ of the
value ranges with $t_i(I_i(A))=I'_{\mu i}(H(A))$ for every object $A$.
For predicate invariants we require
$P_i(A)\Longleftrightarrow P'_{\mu i}(H(A))$.
For the signatures we give a map of axes $\sigma:\Lambda_r\to\Lambda_{r'}$ and
bijections $\beta_\ell:V_\ell\xrightarrow{\sim}V'_{\sigma\ell}$ of the value
ranges, and impose
\begin{equation*}
\ell\in\Lambda_{\mathrm{sel}}
\Longleftrightarrow \sigma\ell\in\Lambda'_{\mathrm{sel}},
\qquad
\beta_\ell(q_\ell(A))=q'_{\sigma\ell}(H(A))
\tag{1.11}\label{eq:1.11}
\end{equation*}
A morphism retains all of these constituent maps.
Even morphisms with the same action on configurations are distinct if their
maps of objects, operations, axes, and so on differ.

\end{definition}

\begin{proposition}[Action of exact core changes]\label{prop:1.29}

A morphism as in Definition~\ref{def:1.28} sends $A_r$ to $A_{r'}$ and
satisfies $H(\mathrm{Obj}_r)\subseteq\mathrm{Obj}_{r'}$.
For every object $A$ and equation index $i$,
\begin{equation*}
E_{r,i}(A)\Longleftrightarrow E_{r',\lambda i}(H(A)),
\qquad
\mathrm{Lawful}_{E_r}(A)\Longleftrightarrow
\mathrm{Lawful}_{E_{r'}}(H(A)).
\tag{1.12}\label{eq:1.12}
\end{equation*}
Moreover, $Q\mapsto e_*Q$ gives a bijection between
$\mathrm{Circ}_{E_r}(A,i)$ and $\mathrm{Circ}_{E_{r'}}(H(A),\lambda i)$.

\end{proposition}

\begin{proof}

Using \eqref{eq:1.7} and \eqref{eq:1.8} in turn gives the equality of base objects.
Letting $H,\Phi$ act on a finite sequence of operations from the base
object and using \eqref{eq:1.2} yields the inclusion of generated object
families.

By the residual comparison in \eqref{eq:1.10} and the preservation and
reflection of zero by the ring isomorphisms, the vanishing of a residual on $W$
is equivalent to its vanishing on $\theta W$.
Since $e$ is a bijection, the quantification over all Atoms corresponds as
well.
Every context $V$ after the change is isomorphic to some $\theta W$, and the
restriction along this isomorphism is a ring isomorphism.
Applying \eqref{eq:1.1} to this isomorphism gives the same equivalence for the
residuals on $V$.
Using the bijection of indices and the preservation of roles yields the
equivalence of lawfulness.

By \eqref{eq:1.8}, membership of Atoms and validity of relations and
identifications are preserved and reflected along $e$.
Hence $\mathrm{Matches}(Q,A)\Longleftrightarrow\mathrm{Matches}(e_*Q,H(A))$ for
both truth signs.
From the equality of codes, the accepted sets satisfy $T'_{\lambda i}=e_*T_i$.
Together with the inverse map of query sequences given by $e^{-1}$, we obtain
the bijection of circuits.

\end{proof}

\begin{definition}[Morphisms of geometries]\label{def:1.30}

A morphism from a geometry $G=(r,\mathcal{R},\mathrm{Ov},k,\mathcal{B})$ to
$G'=(r',\mathcal{R}',\mathrm{Ov}',k',\mathcal{B}')$ consists of a morphism of
cores $h:r\to r'$ and the following data and conditions.
The specified quasi-inverse of $\theta$ is written $\bar\theta$.

\textbf{1. Preservation of coverage requirements.} Each requirement is
preserved forward, using $e$ for the required Atoms,
$(i,a)\mapsto(\lambda i,e(a))$ for the equation coordinates and witnesses, and
$\sigma$ for the signature axes.
The four visibility predicates are also preserved along $W\mapsto\theta W$ and
these maps, and
$B_{\mathcal{R}}(U,W)\Rightarrow B_{\mathcal{R}'}(\theta U,\theta W)$ holds.

\textbf{2. Comparison of overlaps.} For each diagram $U\to W\leftarrow V$ in
the context category after the change, we specify an isomorphism between the
carried overlap chosen before the change and the overlap chosen after the
change; that is,
$\theta(\bar\theta U\times_{\bar\theta W}\bar\theta V)\cong U\times_W V$.
The comparison with the projections uses the counit of the equivalence.

\textbf{3. Change of coefficients.} We specify a ring homomorphism
$\varphi:k\to k'$.

\textbf{4. Agreement of ring presentations.} The raw system after the change
equals the raw system before the change with its coefficients changed by
$\varphi$ and reindexed along $\bar\theta$:
\begin{equation*}
\mathcal{B}'=\bar\theta^{\,*}(\varphi_!\mathcal{B}).
\tag{1.13}\label{eq:1.13}
\end{equation*}
The right-hand side places on a context $V$ after the change the coordinates,
kinds, local data types, and relation indices of $\bar\theta V$, with the
coefficients of the polynomials changed.
The restrictions are built from the restrictions along the morphisms sent by
$\bar\theta$.
The equality covers all of these presentation data.

\textbf{5. Comparison of the data contained in contexts.} To each $W$ we
assign comparison maps for supports, axes, and observables
\[
s_W:\mathrm{Supp}(W)\to\mathrm{Supp}'(\theta W),\quad
a_W:\mathrm{Ax}(W)\to\mathrm{Ax}'(\theta W),\quad
o_W:\mathrm{Obs}(W)\to\mathrm{Obs}'(\theta W)
\]
The readings of the support are preserved with the Atoms carried by $e$, and
the readable elements of axes and observables are preserved as well.
Writing $S_j,A_j$ for the covariant structure maps of supports and axes and
$O_j$ for the contravariant restriction of observables along $j:W'\to W$, we
require
\begin{equation*}
S'_{\theta j}s_{W'}=s_WS_j,\qquad
A'_{\theta j}a_{W'}=a_WA_j,\qquad
O'_{\theta j}o_W=o_{W'}O_j
\tag{1.14}\label{eq:1.14}
\end{equation*}
Since the context category is a preorder category, any two parallel morphisms
that exist are equal.
Hence the commutativity with the projections in item~2 and the unit and
composition coherence of the overlap comparisons are determined uniquely.
The comparison maps for supports and the like are retained as actual maps and
satisfy the conditions of item~5 separately.

\end{definition}

\begin{theorem}[The three-level categories and projections]\label{thm:1.31}

The objects and morphisms of Definitions~\ref{def:1.27}, 1.28, and 1.30 form
categories $B$, $E_{\mathrm{core}}$, and $E_{\mathrm{geom}}$, respectively,
with functors
\begin{equation*}
E_{\mathrm{geom}}
\xrightarrow{\,p\,}
E_{\mathrm{core}}
\xrightarrow{\,q\,}
B
\tag{1.15}\label{eq:1.15}
\end{equation*}
The functor $p$ forgets the choice of geometry and its comparisons and extracts
the core, and $q$ extracts the extraction doctrine, the selected source, and
their maps.

\end{theorem}

\begin{proof}

For $B$ this was checked in Definition~\ref{def:1.27}.
Morphisms of cores are composed by composing the maps of sources, Atoms,
objects, indices, and operations in turn.
For the equivalences of context categories we use the composite equivalence,
and for the ring isomorphisms of each context the ring isomorphisms applied in
turn.
Each of \eqref{eq:1.8}--\eqref{eq:1.11} holds after composition by substituting
the two corresponding equalities.
The bijections of value ranges of functional invariants compose, and so do the
equivalences for predicate invariants.
For identity morphisms we use the respective identity maps.
The unit and associativity laws follow from the composition laws of maps and
the uniqueness of comparison morphisms in a preorder category.

The coefficient maps of morphisms of geometries are composed as ring
homomorphisms.
The comparisons of supports and the like are composed, for example, as
$s''_W=s'_{\theta W}s_W$.
The visibility implications and \eqref{eq:1.14} are closed under this
composition.
For raw systems, reindexing commutes with the change of coefficients, and by
Lemma~\ref{lem:1.26} two successive changes of coefficients equal the change by
the composite homomorphism.
Hence the two instances of \eqref{eq:1.13} yield \eqref{eq:1.13} for the
composite morphism.
The identity and associativity laws hold as well in each component: the
coordinates, the relation polynomials, and the restriction maps.

The two projections extract the corresponding components of these composites.
They are therefore defined on objects and morphisms and preserve identities
and composition.

\end{proof}

The diagram \eqref{eq:1.15} is a construction for treating, in a single
diagram, different object formations with the same extraction and different
local geometries with the same core.
Chapter~\ref{chap:4} studies transport, which builds the objects and morphisms
of the upper levels from a change of the base, through its existence and
universal property.

\section{Model synchronization: the semantics of reads and updates}\label{sec:1.8}

In \S\ref{sec:1.7} we compared choices of reading, such as extraction and
equations.
From here on we compare states and operations under semantics defined
independently of AAT.
We take up lenses, which treat the consistency of reads and updates, and
protocols, which treat executions chaining named operations together with their
observations.

In \S\ref{sec:1.10} we construct from both of them objects, operations, and
Laws organized by role, and show that the constructed Laws hold if and only if
the original laws hold.
We further check that typed morphisms between the constructed objects are in
one-to-one correspondence with the original semantics-preserving morphisms.
Through this correspondence, the results obtained in later chapters from the
common theorems, such as the classification of changes and local
reconstruction, are brought back to the original semantics.

For lenses, we first determine a decomposition of states and the maps
preserving reads and updates.
In model synchronization, part of the internal state is displayed, and changes
to the displayed value are reflected back into the internal state.
The displayed part is called the view.
When the consistency of reads and updates is specified by the three lens laws,
a state decomposes into a displayed value and a component preserved by updates.

\begin{definition}[Total lenses and their morphisms]\label{def:1.32}

Fix a set $V$ of views and a reference value $v_0\in V$.
A total lens is a state set $C$, a read $g:C\to V$, and an update
$p:C\times V\to C$ satisfying
\begin{equation*}
p(c,g(c))=c,\qquad
g(p(c,v))=v,\qquad
p(p(c,v),w)=p(c,w)
\tag{1.16}\label{eq:1.16}
\end{equation*}
for all $c,v,w$.
The three equations express, in order, the following.

\begin{enumerate}
\item Writing back the current displayed value unchanged does not change the
state.
\item Reading after an update yields the specified displayed value.
\item Updating twice in succession gives the same result as updating once with
the last displayed value.
\end{enumerate}

In standard lens terminology, this is a total, very well-behaved lens
(\cite[\href{https://www.cis.upenn.edu/~bcpierce/papers/newlenses-full.pdf}{\S3.1}]{FGMPS04}).
For the lenses treated below, the reference fiber
$K_L=\{c\in C\mid g(c)=v_0\}$ is required to be finite.
This is the set of states remaining when the display is fixed at the reference
value.
The set $V$ of displayed values itself may be infinite.
The finiteness of the reference fiber is used later to describe
semantics-preserving morphisms by finite tables.

A morphism $h:L\to L'$ over the fixed view is a map of states $h:C\to C'$
satisfying
\begin{equation*}
g'h=g,\qquad h(p(c,v))=p'(h(c),v)
\tag{1.17}\label{eq:1.17}
\end{equation*}
Identities and composites are the identities and composites of state maps, and
by substitution in \eqref{eq:1.17} they satisfy \eqref{eq:1.17} again.
This category is written $\mathrm{Lens}(V,v_0)$.

\end{definition}

\begin{proposition}[Product decomposition of lenses]\label{prop:1.33}

For a total lens $L=(C,g,p)$, the maps
\begin{equation*}
\begin{aligned}
\eta_L:C&\longrightarrow V\times K_L,
&c&\longmapsto(g(c),p(c,v_0)),\\
\zeta_L:V\times K_L&\longrightarrow C,
&(v,k)&\longmapsto p(k,v)
\end{aligned}
\tag{1.18}\label{eq:1.18}
\end{equation*}
are mutually inverse bijections. In this decomposition the read is the first
projection, and the update is $((v,k),w)\mapsto(w,k)$.

\end{proposition}

\begin{proof}

Since $g(p(c,v_0))=v_0$, the map $\eta_L$ is well defined.
The three laws give
\[
p(p(c,v_0),g(c))=p(c,g(c))=c
\]
Hence $\zeta_L\eta_L=\mathrm{id}$. On the other hand, for $k\in K_L$ we have
$g(p(k,v))=v$ and
$p(p(k,v),v_0)=p(k,v_0)=p(k,g(k))=k$, so
$\eta_L\zeta_L=\mathrm{id}$ holds as well.
The formula for the read follows from the first component. After an update the
second component $p(p(c,w),v_0)=p(c,v_0)$ is unchanged, and the first component
becomes $w$.

\end{proof}

The second component of the product decomposition represents the information
retained when the displayed value changes.
A map preserving reads and updates is likewise determined solely by how it acts
on this component.

\begin{proposition}[Recovery of morphisms from the reference fiber]\label{prop:1.34}

For any two lenses, the restriction
\[
\mathrm{res}:\mathrm{Hom}(L,L')
\longrightarrow\mathrm{Map}(K_L,K_{L'}),\qquad
h\longmapsto h|_{K_L}
\]
is a bijection. Here $\mathrm{Map}(K,K')$ denotes the set of all maps between
the sets.
The inverse is given, for any $t:K_L\to K_{L'}$, by
\begin{equation*}
\mathrm{ext}(t)(c)=p'\bigl(t(p(c,v_0)),g(c)\bigr)
\tag{1.19}\label{eq:1.19}
\end{equation*}
This is the operation of updating once to the reference view, applying $t$ to
that state, and then updating back to the original displayed value.

\end{proposition}

\begin{proof}

A morphism $h$ preserves reads and therefore sends the reference fiber into the
reference fiber.
In the product decomposition of Proposition~\ref{prop:1.33}, $\mathrm{ext}(t)$
is $(v,k)\mapsto(v,t(k))$, which preserves reads and updates.
Hence \eqref{eq:1.19} indeed gives a morphism of lenses.

For $k\in K_L$ we have $p(k,v_0)=k$ and
$p'(t(k),v_0)=t(k)$, so $\mathrm{res}(\mathrm{ext}(t))=t$.
In the other direction, \eqref{eq:1.17} and the three laws give
\[
\begin{aligned}
\mathrm{ext}(\mathrm{res}(h))(c)
&=p'(h(p(c,v_0)),g(c))\\
&=p'(p'(h(c),v_0),g(c))\\
&=p'(h(c),g(c))
=p'(h(c),g'(h(c)))=h(c)
\end{aligned}
\]
That both restriction and extension preserve identities and composition is seen
from the product decomposition.

\end{proof}

As a normal form including visible changes, we treat here self-changes of the
same product lens $V\times K$ for a finite set $K$.
When $u:V\xrightarrow{\sim}V$ and $h:V\times K\xrightarrow{\sim}V\times K$
satisfy
\begin{equation*}
gh=ug,\qquad h(p(c,v))=p(h(c),u(v))
\tag{1.20}\label{eq:1.20}
\end{equation*}
there is a unique permutation $\phi:K\xrightarrow{\sim}K$ with
\[
h(v,k)=(u(v),\phi(k))
\]
Preservation of updates gives that the permutation of the hidden component
obtained over each view equals the permutation over the reference view.
Reading the second component at $v=v_0$ also gives the uniqueness of that
permutation.

\begin{example}[Changes preserving only reads]\label{ex:1.35}

Consider the product lens with $V=K=\{0,1\}$, where $\oplus$ denotes addition
modulo $2$.
The bijection $h(v,k)=(v,k\oplus v)$ preserves reads, but
\[
h(p((0,0),1))=(1,1),\qquad
p(h(0,0),1)=(1,0)
\]
The bijections over the fixed view preserving reads number $2\cdot2=4$,
choosing the identity or the swap independently on the two fibers.
Those preserving updates as well number $2$, making the same choice on both
fibers.
By tying the different fibers together, the update constrains the permitted
changes.

\end{example}

\section{Protocols: named operations and all executions}\label{sec:1.9}

A protocol has individual operations and execution sequences chaining them.
When cycles can be repeated, finitely many operations give rise to execution
sequences of arbitrary length.
The meaning of all executions is fixed by specifying the action of each
operation and the pairs of paths to be treated as equal executions.
For this description we use the standard constructions of presenting a category
by a graph with path equations and of realization by set-valued functors
(\cite[\href{https://arxiv.org/pdf/1009.1166v3}{\S\S3.2, 3.4--3.5}]{Spivak12}).

\begin{definition}[Realizations of protocols]\label{def:1.36}

A finite directed multigraph $Q$ is given by a vertex set $Q_0$, a set $Q_1$ of
named edges, and start and end maps.
A path is a composable finite sequence of edges, and the empty sequence at each
vertex is the identity path.
We specify finitely many pairs of paths
$\Pi=\{(\ell_j,r_j)\}_{j\in J_\Pi}$ with the same start and end.

On all paths we impose the smallest equivalence relation that identifies the
specified pairs and is preserved by composition on either side.
The category with the vertices as objects and the equivalence classes of paths
as morphisms is called $\mathcal{C}_Q$.
Composition is induced by concatenation of paths and is well defined on
equivalence classes by congruence.
Associativity and the unit laws follow from concatenation of paths.

Fix a functor $O:\mathcal{C}_Q\to{\mathbf{Set}}$ of observation targets.
A realization of the protocol is a pair of a functor
$X:\mathcal{C}_Q\to{\mathbf{Set}}$ whose state set $X(v)$ at each vertex is
finite and a natural transformation $o_X:X\Rightarrow O$.
For an edge $e:v\to w$, the map $X(e)$ represents the transition of states, and
$o_X(v)$ gives the observed value readable from a state at the vertex $v$.
Naturality of the observation is the condition that transitioning the state
and then observing agrees with carrying the observed value by $O(e)$.

A morphism $a:(X,o_X)\to(Y,o_Y)$ is a natural transformation $a:X\Rightarrow Y$
with $o_Ya=o_X$. Its components $a_v$ are general maps of states.
Naturality preserves executions, and the latter equation preserves
observations.
With componentwise identities and composition we obtain the category
$\mathrm{Prot}(Q,\Pi,O)$.

\end{definition}

\begin{proposition}[Realization from finite operation tables]\label{prop:1.37}

Assign to each vertex a finite set $S_v$, to each edge $e:v\to w$ a map
$T_e:S_v\to S_w$, and to each vertex an observation $b_v:S_v\to O(v)$.
Define the action of paths by composing the edge maps, and assume the
following.
\begin{equation*}
T_{\ell_j}=T_{r_j}\quad(j\in J_\Pi),\qquad
b_wT_e=O(e)b_v\quad(e:v\to w).
\tag{1.21}\label{eq:1.21}
\end{equation*}
Then there is exactly one realization with these states, edges, and
observations.

\end{proposition}

\begin{proof}

Assign the identity map to the empty path and $T_{e_n}\cdots T_{e_1}$ to
$e_n\cdots e_1$.
This is a functor from the free path category to the category of sets.
The pairs of paths inducing the same map form an equivalence relation preserved
by composition on either side and containing the declared relations of
\eqref{eq:1.21}.
The action is therefore well defined on the equivalence classes of paths of
$\mathcal{C}_Q$.

Commutativity of the observations also follows by induction on the length of
the path.
Indeed, for $p:v\to w$ and $e:w\to z$, if $b_wT_p=O(p)b_v$, then
$b_zT_eT_p=O(e)b_wT_p=O(e)O(p)b_v$.
Hence $b$ is a natural transformation.
Since every path is a composite of edges, uniqueness of the extension follows
as well.

\end{proof}

\begin{proposition}[Recovery of semantics-preserving morphisms from vertex maps]\label{prop:1.38}

For realizations $X,Y$, if vertexwise maps $a_v:X(v)\to Y(v)$ satisfy
\begin{equation*}
a_wX(e)=Y(e)a_v,\qquad o_Y(v)a_v=o_X(v)
\tag{1.22}\label{eq:1.22}
\end{equation*}
for all named edges $e:v\to w$ and all vertices, they extend to a unique
morphism of protocols.

\end{proposition}

\begin{proof}

Commutativity for the empty path is a property of identity maps, and for a
path $p:v\to w$ and an edge $e:w\to z$ we have
\[
a_zX(e)X(p)=Y(e)a_wX(p)=Y(e)Y(p)a_v
\]
By induction we obtain naturality with respect to all paths.
Since the actions of both realizations depend on the equivalence class of
a path, naturality holds for all morphisms of the quotient category.
Preservation of observations is exactly \eqref{eq:1.22}.
A natural transformation is determined by its components at the objects, so it
is unique.

\end{proof}

Since the vertices, the named edges, the declared relations, and the state sets
are all finite, the inputs of Propositions~\ref{prop:1.37} and 1.38 are given
by finitely many tables.
On the other hand, paths repeating cycles can have arbitrary length.
Conclusions about all executions are obtained not by enumerating them but from
path induction and the congruence relation.

We regard a semantics-preserving morphism connecting two realizations as an
adapter.
When comparing adapters $q:X\to Y$ and $q':X'\to Y'$, the condition for
changes $a:X\to X'$ and $b:Y\to Y'$ to preserve the adapters is
\begin{equation*}
bq=q'a.
\tag{1.23}\label{eq:1.23}
\end{equation*}
That is, applying the adapter and then the change agrees with applying the
change and then the adapter after the change.
As an equation of natural transformations, this is equivalent to
$b_vq_v=q'_va_v$ at each vertex.
By \eqref{eq:1.22}, both sides commute with all executions.
When studying invertible changes we restrict $a,b$ to isomorphisms.
The adapters $q,q'$ themselves continue to be treated as general
semantics-preserving morphisms.

\section{Constructing Atoms, operations, and Laws from semantics}\label{sec:1.10}

To express lenses and protocols as objects of AAT, one must decide where to
keep the operation names, the state values, the state sets, and the meaning of
the operations.
Here we place the type roles and the operation names in Atoms, and the
individual state values in the source.
The state sets are kept in architecture objects, and the functions such as
reads, updates, and transitions in operations.
This arrangement keeps the vocabulary finite while still handling general maps
sending states to other states.

\begin{construction}[Type roles and named operations]\label{cons:1.39}

We record the difference of object roles in the target of the relations, and
the difference of operation names in the image of a special Atom.
This assigns a morphism of configurations to each operation while still
distinguishing operations that act as the same function on states.

The Atom set of a lens $L=(C,g,p)$ consists of the seven distinct symbols
\[
\mathrm{At}_L=
\{*,\mathrm{state},\mathrm{view},\mathrm{read},
\mathrm{write},\mathrm{get},\mathrm{put}\}
\]
Taking this set as the value range of all five coordinates and the coordinate
maps to be the identity satisfies Definition~\ref{def:1.1}.
The vocabulary of a protocol is built in the same way from the distinct
symbols
\[
\mathrm{At}_Q=
\{*\}\sqcup
\{\mathrm{state}_v,\mathrm{observation}_v,\mathrm{observe}_v
\mid v\in Q_0\}
\sqcup\{\mathrm{edge}_e\mid e\in Q_1\}
\]
Both are nonempty finite sets.

For a selected role $t\in\mathrm{At}$, define the configuration
\[
C_t=(\mathrm{At},R_t,\varnothing),\qquad
R_t(a,b)\ \Longleftrightarrow\ b=t
\]
All Atoms are contained in the family, and the relations point to the selected
role.
Which role a configuration represents can therefore be seen from the target
$t$ of its relations.
The set corresponding to the role is placed in the structure data $S_{A_t}$,
with $Q_{A_t}=v_0$ for lenses and $Q_{A_t}=*$ for protocols.
The sets for the roles are as follows.

\begin{xltabular}{\linewidth}{@{}LLL@{}}
\toprule
Semantics & Role & Set kept \\
\midrule
\endhead
lens & $\mathrm{state}$ & The states $C$ \\
\addlinespace[3pt]
lens & $\mathrm{view}$ & The displayed values $V$ \\
\addlinespace[3pt]
lens & $\mathrm{read}$ & The inputs $C$ of the read \\
\addlinespace[3pt]
lens & $\mathrm{write}$ & The inputs $C\times V$ of the update \\
\addlinespace[3pt]
protocol & $\mathrm{state}_v$ & The states $X(v)$ at the vertex $v$ \\
\addlinespace[3pt]
protocol & $\mathrm{observation}_v$ & The observed values $O(v)$ at the vertex $v$ \\
\bottomrule
\end{xltabular}

At state and read of the lens we place the same set $C$; it serves as the role
of update results and as the role of read inputs, respectively.
At write we place the pairs of a current state and a specified displayed value.

Between these objects we provide the following named operations.

\begin{xltabular}{\linewidth}{@{}LLL@{}}
\toprule
Operation name & Roles of domain and codomain & Function giving the meaning \\
\midrule
\endhead
$\mathrm{get}$ & $\mathrm{read}\to\mathrm{view}$ & $g$ \\
\addlinespace[3pt]
$\mathrm{put}$ & $\mathrm{write}\to\mathrm{state}$ & $(c,v)\mapsto p(c,v)$ \\
\addlinespace[3pt]
$\mathrm{edge}_e$ ($e:v\to w$) & $\mathrm{state}_v\to\mathrm{state}_w$ & $X(e)$ \\
\addlinespace[3pt]
$\mathrm{observe}_v$ & $\mathrm{state}_v\to\mathrm{observation}_v$ & $o_X(v)$ \\
\bottomrule
\end{xltabular}

To an operation $s\to t$ with name $m$ we assign the map on Atoms
\begin{equation*}
d_m(a)=
\begin{cases}
m & a=*,\\
t & a\ne *
\end{cases}
\tag{1.24}\label{eq:1.24}
\end{equation*}
Every domain role $s$ in the table differs from $*$.
Hence $R_s(a,b)$ implies $b=s$, so $d_m(b)=t$, that is,
$R_t(d_m(a),d_m(b))$.
Preservation of the family follows because it contains all Atoms, and
preservation of the identifications because they are empty.
Thus $d_m:C_s\to C_t$ is indeed a morphism of configurations.

We take an operation to be the tuple of its name, the objects at its two ends,
$d_m$, and the function in the table giving its meaning.
This provides the operation reading of Definition~\ref{def:1.7}.
Since $d_m(*)=m$, distinct names are distinguished even as morphisms of
configurations.

\end{construction}

\begin{construction}[Extraction sources carrying state values]\label{cons:1.40}

Take the sources of a lens and of a protocol to be
\begin{equation*}
\begin{aligned}
\mathrm{Src}_L
&=1\sqcup C\sqcup V\sqcup(C\times V),\\
\mathrm{Src}_X
&=1\sqcup\bigsqcup_{v\in Q_0}X(v)
 \sqcup\bigsqcup_{e\in Q_1}X(\mathrm{src}(e))
\end{aligned}
\tag{1.25}\label{eq:1.25}
\end{equation*}
Each summand carries a tag: point, state, view, update input, or edge input.
Normalization is the identity, and the vocabulary, meaning, and resolution
parameters are one-element sets.
The predicates $V_D,M_D,R_D$ of Definition~\ref{def:1.3} are always true, and
$E_D$ is defined by the following table.

\begin{xltabular}{\linewidth}{@{}LL@{}}
\toprule
Tag of the source & Atoms extracted \\
\midrule
\endhead
Point (both semantics) & All Atoms of the vocabulary \\
\addlinespace[3pt]
State $c$ of the lens & $\mathrm{state},\mathrm{read},\mathrm{get}$ \\
\addlinespace[3pt]
View $v$ of the lens & $\mathrm{view}$ \\
\addlinespace[3pt]
Update input $(c,v)$ of the lens & $\mathrm{write},\mathrm{put}$ \\
\addlinespace[3pt]
State $c$ at a vertex $v$ of the protocol & $\mathrm{state}_v,\mathrm{observation}_v,\mathrm{observe}_v$ \\
\addlinespace[3pt]
Input state $c$ of an edge $e$ of the protocol & $\mathrm{edge}_e$ \\
\bottomrule
\end{xltabular}

Taking the point as the selected source, the extracted family is the whole
vocabulary, which is finite.
The state values are kept in the sources \eqref{eq:1.25}, and the extraction
predicate reads their tags.

A morphism of lenses $h:L\to L'$ determines a map of sources using the
identity, $h$, the identity, and $h\times\mathrm{id}_V$ on the respective
summands.
A morphism of protocols $a:X\to Y$ uses $a_v$ on the state summands and
$a_{\mathrm{src}(e)}$ on the edge input summands, and sends the point to the
point.
Together with the identity map on Atoms, these give morphisms as in
Definition~\ref{def:1.27}.
Commutation with normalization is a property of identity maps, and the
equivalence of extraction follows because the tags are preserved.

Identities and composites of source maps agree with identities and composites
on each summand.
Moreover, reading the state summands off the source map recovers $h$ or all
the $a_v$.
The correspondence is therefore faithful on morphisms, and it also retains
non-injective semantics-preserving maps.

\end{construction}

\begin{proposition}[Laws expressing the laws of the semantics]\label{prop:1.41}

Fix the sets $C,V$ of a lens or the family of state sets $(S_v)$ of a protocol.
Consider operation data $d$ on which no laws have yet been imposed:
$d=(g,p)$ for a lens and $d=((T_e),(b_v))$ for a protocol.
From these data one can construct an architecture object $A_d$ and an equation
system such that being Lawful is equivalent to the laws of the respective
semantics.

\end{proposition}

\begin{proof}[Proof and construction]

We split each law into equations indexed by input values and build residuals
that are zero when the equation holds and one when it does not.

Let the configuration of $A_d$ be $C_*$ and its structure data the operation
data $d$.
As the quantity we keep the reference view for a lens and the element of a
one-element set for a protocol.
The indices of the index set $K$ of equations are the pairs of one of the
following equations and an input value.
\begin{equation*}
\begin{array}{ll}
\text{lens:}&
p(c,g(c))=c,\quad
g(p(c,v))=v,\quad
p(p(c,v),w)=p(c,w),\\[2pt]
\text{protocol:}&
T_{\ell_j}(c)=T_{r_j}(c),\quad
b_w(T_e(c))=O(e)(b_v(c)).
\end{array}
\tag{1.26}\label{eq:1.26}
\end{equation*}
For the lens, $c$, $(c,v)$, and $(c,v,w)$ range over the inputs, respectively.
For the protocol, the declared relations together with a state at their start,
and the edges $e:v\to w$ together with $c\in S_v$, range likewise.
All indices are required.

There is a single context $W$, whose support reads all Atoms from a
one-element set.
The axes and observables are also one-element sets.
This is a context of $A_d$, whose family contains all Atoms.
The constant ring presheaf and the equation coordinates are given by
\[
O_E(W)=\mathbb{Z}[x_{(i,a)}\mid (i,a)\in K\times\mathrm{At}],
\qquad \nu_{W,i,a}=x_{(i,a)}
\]
When an object $A$ under evaluation carries operation data of the same type,
let $P_i(A)$ be the equation of index $i$ for those data.
For objects without structure data of the corresponding type, $P_i(A)$ is
defined to be false.
The residuals are the constant polynomials
\begin{equation*}
\varepsilon_{W,A,i,a}=
\begin{cases}
0 & P_i(A),\\
1 & \text{otherwise}
\end{cases}
\tag{1.27}\label{eq:1.27}
\end{equation*}
Since the restrictions are identities, \eqref{eq:1.1} holds.

If all the equations of \eqref{eq:1.26} hold, all residuals of $A_d$ are zero.
Conversely, if lawfulness holds, we evaluate each $i$ at the context $W$ and
the Atom $*$.
Since $0\ne1$ in the polynomial ring over $\mathbb{Z}$, the vanishing of
\eqref{eq:1.27} implies $P_i(A_d)$.
This yields all the equations of \eqref{eq:1.26}.
In the protocol case, by Proposition~\ref{prop:1.37} these agree with the
conditions defining a functor and a natural observation.

\end{proof}

For example, if $C=V=\{0,1\}$, $g(c)=c$, and $p(c,v)=c$, then
$g(p(0,1))=0\ne1$. The residual corresponding to this input is indeed $1$, and
the constructed object is not Lawful.

\begin{definition}[Morphisms between the constructed typed objects]\label{def:1.42}

We write $\mathcal{T}(L)$ or $\mathcal{T}(X)$ for the totality of the roles,
objects, and named operations with their meanings from
Construction~\ref{cons:1.39}.
A typed morphism between them is a family of maps on the sets of the roles
satisfying the following conditions.

\begin{itemize}
\item \textbf{Lens}: the maps at state and read are the same $h:C\to C'$, the
map at view is the identity, and the map at write is $h\times\mathrm{id}_V$.
We require the two commutativities for get and put.
\item \textbf{Protocol}: the maps at state are $a_v:X(v)\to Y(v)$ and the maps
at observation are the identity, and we require the commutativity for each
edge and observe.
\end{itemize}

With componentwise identities and composition, these objects and typed
morphisms form a category.

Construction~\ref{cons:1.40} connects the state maps given by these typed
morphisms with the source maps of the extraction doctrines.

\end{definition}

\begin{proposition}[Recovery of semantics-preserving morphisms]\label{prop:1.43}

Reading the sets of the roles and the meanings of the operations off the
constructed objects recovers the original state sets, reads, updates, edge
actions, and observations. Moreover, the following natural bijections hold.
\begin{equation*}
\begin{aligned}
\mathrm{Hom}_{\mathrm{Lens}(V,v_0)}(L,L')
&\simeq\mathrm{Hom}_{\mathrm{typ}}(\mathcal{T}(L),\mathcal{T}(L')),\\
\mathrm{Hom}_{\mathrm{Prot}(Q,\Pi,O)}(X,Y)
&\simeq\mathrm{Hom}_{\mathrm{typ}}(\mathcal{T}(X),\mathcal{T}(Y)).
\end{aligned}
\tag{1.28}\label{eq:1.28}
\end{equation*}

\end{proposition}

\begin{proof}

For the objects, it suffices to project the sets and functions encoded in
Construction~\ref{cons:1.39}.
The operation names are also recovered from the value of \eqref{eq:1.24} at
$*$.

A semantics-preserving morphism $h$ of lenses determines the maps on the four
roles of Definition~\ref{def:1.42}.
The commutativities for get and put are \eqref{eq:1.17}, so we obtain a typed
morphism.
Conversely, taking the map at state from a typed morphism gives a morphism of
lenses by the two commutativities.
The maps on the other roles are determined by the map at state, so the two
operations are mutually inverse.

For protocols as well, taking the vertex components of a semantics-preserving
morphism as the maps at state gives a typed morphism.
Conversely, the commutativities for edge and observe of a typed morphism are
\eqref{eq:1.22}, and by Proposition~\ref{prop:1.38} a natural transformation
with respect to all paths is recovered uniquely.
These two constructions are also mutually inverse at the components of each
vertex.

Both send identities to identities and composites to componentwise composites.
Hence \eqref{eq:1.28} is natural with respect to pre- and postcomposition at
both ends.

\end{proof}

A typed morphism of lenses can be recovered, by Proposition~\ref{prop:1.34},
from the map between reference fibers alone.
A typed morphism of protocols extends to all executions from finitely many
vertex maps and the commutativities for the generating edges.
For a lens with a visible change $u$, taking the map at view to be $u$ and the
map at write to be $h\times u$ turns the commutativities for the two
operations into \eqref{eq:1.20}.
Commutativity with protocol adapters likewise agrees with \eqref{eq:1.23} upon
reading back the vertex components.

\section*{Summary of the Chapter}

In this chapter we made explicit the choices of extraction, object formation,
operations, Laws, and locality, and constructed relative architecture and its
comparisons. The main results are the following four.

\begin{itemize}
\item \textbf{Generation of objects.} From an extracted finite family of Atoms
we form the base object. The smallest family of objects containing it and
closed under the permitted operations equals the totality of objects reachable
by finite sequences of operations (Theorem~\ref{thm:1.18}).
Under soundness of the detector and completeness for the required equations,
the validity of the Laws is equivalent to the absence of required circuits
(Proposition~\ref{prop:1.16}).
\item \textbf{Construction of local structure.} From contexts and covers we
constructed a site, and from restriction maps compatible with the coordinates
and structural relations a presheaf of rings. We also constructed the
sheafification, which provides the gluing of local data
(\S\S\ref{sec:1.5}--\ref{sec:1.6}).
\item \textbf{Comparison of readings.} We defined objects and morphisms at
each of the stages of extraction, core, and geometry, and constructed the
projection functors connecting them (Theorem~\ref{thm:1.31}).
\item \textbf{Recovery of the semantics of CS.} We described the lenses and
protocols treated in this chapter in the common framework of objects,
operations, and Laws. The constructed Laws hold if and only if the original
laws hold (Proposition~\ref{prop:1.41}), and typed morphisms correspond
one-to-one to the original semantics-preserving morphisms
(Proposition~\ref{prop:1.43}).
\end{itemize}

In the next chapter we construct an ideal from the symbolic equation
coordinates.
Under a realization condition connecting the evaluation of the coordinates
with the residuals of each object, we describe the part satisfying the Laws
inside the local geometry.
